% chapters/ch17_mock_interviews.tex — Chapter 17: 5 Full Mock Interviews
% Scope (PLAN.md): screening call, technical round 1 (C + protocols), technical
% round 2 (domain + projects), coding round, managerial round. Each: ~45-min
% question script, scoring rubric, full answer key (cross-referenced to the
% chapters), common-mistake notes. Budget 20 pp. No new code (references Ch13).
\chapter{Five Full Mock Interviews}
\label{ch:mock-interviews}

\begin{partnote}
\textbf{Part E --- Interview Execution.} Five simulated interviews to run on
paper (or with a partner), ~45 minutes each: \textbf{(1)} screening call,
\textbf{(2)} technical round 1 --- C \& protocols, \textbf{(3)} technical round 2
--- domain \& projects, \textbf{(4)} coding round, \textbf{(5)} managerial round.
Each has a \emph{question script}, a \emph{scoring rubric}, an \emph{answer key}
(cross-referenced to the chapters), and \emph{common-mistake} notes.

\medskip
\textbf{Rules:} answer \emph{aloud}, time yourself, and score honestly against
the rubric before reading the key. Project- and behavioural-specific answers
(the ``tell me about your DFCC work'' depth) live in Chapters 15--16, which are
built only from \code{inputs/INTAKE.md}; this chapter rehearses the
\emph{technical} rounds and the \emph{structure} of the behavioural one. Target:
$\ge$80\% on each rubric before you sit the real thing.
\end{partnote}

%=====================================================================
\section{Mock 1 --- Screening Call (Recruiter + Technical Phone)}
%=====================================================================
\emph{~30--45 min, often a lead engineer. Goal: confirm the resume is real,
filter obvious gaps, check communication. Answers must be crisp and
spoken-length.}

\subsection*{Question script}
\question{S1}{Walk me through your background in two minutes.}
A tight narrative: current role (TASL UAV R\&D --- DFCC/CCDL/SPIL validation,
RS422 Embedded C), prior role (BEL LCA-Tejas DFCC V\&V, ARM9 firmware, the
Python automation), core strengths (embedded C, protocols, board bring-up,
DO-178B V\&V), and what you're looking for. Practise this until it's 90 seconds,
not five minutes.

\question{S2}{What does your day-to-day actually involve?}
Concrete activities: board bring-up and flashing (ST-Link and J-Link),
RS422/CAN interface validation, debugging with scope, DMM, and logic analyser,
packet-level analysis. Specifics signal a real engineer.

\question{S3}{What is \code{volatile} and when do you use it? (warm-up
technical)}
Visibility, not atomicity; memory-mapped registers, ISR-shared globals
(Chapter 1/4). Two sentences, confident.

\question{S4}{Why is RS422 used on flight computers rather than RS232?}
Differential noise immunity, long range, full-duplex (Chapter 5). One paragraph.

\question{S5}{What's your experience with DO-178B?}
Level A verification, MC/DC coverage, requirements traceability, ATP/QTP
(Chapter 12) --- stated at the level the resume claims, no more.

\question{S6}{Why are you leaving / looking? \quad Notice period? \quad Salary
expectation?}
\textbf{[From INTAKE.md / Chapter 16 --- have crisp, honest, non-negative
answers ready. Do not improvise these on the call.]}

\begin{scorecard}
\textbf{Rubric (10 pts):} 2-min pitch is $\le$2 min and structured (2);
day-to-day is concrete, not vague (2); \code{volatile} correct incl.\ ``not
atomic'' (1); RS422 rationale correct (1); DO-178B claim matches resume, not
over-stated (2); logistics answers ready and positive (2). \textbf{Pass
$\ge$8.}
\end{scorecard}

\textbf{Common mistakes:} rambling past two minutes; vague ``I do testing''
(say \emph{what} you test and \emph{how}); over-claiming DO-178B depth you
can't defend; negative reasons for leaving; not knowing your own notice period/
number. The screen is filtering for \emph{red flags}, so be concise, concrete,
and positive.

%=====================================================================
\section{Mock 2 --- Technical Round 1 (Embedded C \& Protocols)}
%=====================================================================
\emph{~45 min, panel of engineers. The make-or-break round. Whiteboard-level C
and protocol depth.}

\subsection*{Question script}
\question{T1.1}{Explain \code{volatile} vs atomic. An ISR sets a flag the main
loop polls --- what can go wrong?}
Visibility vs atomicity vs ordering; torn reads, reordering; fix with barrier/
critical section (Chapters 1, 4). \emph{This is the single highest-value
answer in the book --- nail it.}

\question{T1.2}{What is structure padding? Given \code{uint8\_t; uint32\_t;
uint16\_t}, the size and how to shrink it.}
12 bytes $\rightarrow$ 8 by reordering largest-first (Chapter 1).

\question{T1.3}{Read these declarations: \code{const char *p} vs
\code{char *const p}.}
Right-to-left (Chapter 1).

\question{T1.4}{Describe a UART frame and how you'd compute the baud divider /
check the error.}
Start/data-LSB-first/parity/stop; \code{BRR=fck/baud}, error in ppm, $<$2\%
rule (Chapter 5).

\question{T1.5}{How does CAN arbitration work? Which ID wins and why no
collision?}
Dominant beats recessive, bitwise on the ID, loser backs off (Chapter 7).

\question{T1.6}{When would you choose SPI vs I2C vs UART?}
Clocking model, throughput, pins, device count (Chapter 6).

\question{T1.7}{A node keeps going bus-off on CAN. How do you diagnose it?}
Bit-timing/sample-point mismatch, termination, transceiver --- scope CAN\_H/L,
check bit timing (Chapter 7).

\question{T1.8}{What's the difference between an MPU and an MMU? Which does your
LPC3250 have?}
Protection vs translation; ARM9 has an MMU (Chapter 3).

\question{T1.9}{Why does telemetry use UDP, not TCP?}
Freshness over reliability; no retransmit stalls (Chapter 8).

\question{T1.10}{Why can't you cast a received 1553/ARINC buffer onto a C
struct?}
Endianness + implementation-defined bitfields + wire quirks; parse with
shifts/masks (Chapters 1, 8).

\begin{scorecard}
\textbf{Rubric (20 pts, 2 each):} T1.1 must separate visibility/atomicity/
ordering \emph{and} name a fix (the keystone); each of T1.2--T1.10 correct and
confident. \textbf{Pass $\ge$16; T1.1 is mandatory --- a miss there caps the
round.}
\end{scorecard}

\textbf{Common mistakes:} ``\code{volatile} makes it thread-safe'' (instant
red flag); forgetting the \emph{tail} padding in T1.2; saying sub-priority
pre-empts on CAN/NVIC; ``just lower the baud/clock'' as a first reflex instead
of measuring; trusting a struct overlay on wire data. Reason aloud and reach for
the scope, not guesses.

%=====================================================================
\section{Mock 3 --- Technical Round 2 (Domain \& Projects)}
%=====================================================================
\emph{~45 min, senior/principal engineer. Depth on systems, real-time, control,
and \emph{your} projects. This is where honesty about scope scores.}

\subsection*{Question script}
\question{T2.1}{Explain priority inversion and how it's fixed. Have you seen it?}
High blocked on low's mutex, medium preempts low; priority inheritance; Mars
Pathfinder (Chapter 9). \textbf{[Tie to your real timing-defect work via
INTAKE.md.]}

\question{T2.2}{Walk me through your DPRAM architecture and how it hit sub-1\,ms
latency.}
Single-slot flag handshake, ordering + barrier, poll period bounds latency
(Chapter 9). \textbf{[Project specifics from INTAKE.md / Chapter 15.]}

\question{T2.3}{What is CCDL and what exactly do you verify on it?}
Cross-channel exchange + voting each frame; integrity, timing, failover;
median/threshold vote, fail-operational vs fail-safe (Chapters 8, 10).

\question{T2.4}{Explain a PID loop. What does each term cost, and what's integral
windup?}
P/I/D intuition, the cost of each term, anti-windup (Chapter 10).

\question{T2.5}{How would you verify ``fail over to backup within 50\,ms''?}
Fault-injection HIL measuring the time + MC/DC on detection + worst-case timing
analysis (Chapter 12).

\question{T2.6}{How does your Python suite work, and how did it cut regression
60\%?}
\code{struct} framing, CRC matched to firmware, data-driven cases, unattended
repeatable runs (Chapter 14). \textbf{[Real numbers from INTAKE.md.]}

\question{T2.7}{Where does your knowledge end on the control law itself?}
\textbf{Honest boundary:} ``I verify the control-law outputs and the loop
behaviour; the law \emph{design} is the systems team's, and it's typically
state-space/gain-scheduled --- but I can explain how I'd verify its response.''
(Chapter 10, Q3.3.)

\begin{scorecard}
\textbf{Rubric (14 pts, 2 each):} correct mechanism for T2.1--T2.5; project
answers are specific and consistent with the resume (T2.2, T2.6); T2.7 shows a
\emph{clean honest boundary} rather than bluffing. \textbf{Pass $\ge$11. A smooth
``here's where my knowledge ends'' scores higher than a confident wrong
answer.}
\end{scorecard}

\textbf{Common mistakes:} bluffing past your real scope (a senior interviewer
will dig and expose it); project claims that don't match the resume; describing
the DPRAM handshake without the \emph{ordering/barrier} that makes it safe;
``more integral is better''; treating redundancy as ``always keeps working''
(dual = fail-safe, triple = fail-operational). Knowing the boundary is the
senior signal.

%=====================================================================
\section{Mock 4 --- Coding Round}
%=====================================================================
\emph{~45 min, whiteboard or shared editor. Write compilable C, talk through
edge cases. All of these are in Chapter 13 --- here they're timed under
pressure.}

\subsection*{Question script (write and dry-run each)}
\question{C1}{Implement a lock-free SPSC ring buffer for a UART RX ISR. Why no
lock?}
Single writer per index (Chapter 13 \#31 / Chapter 4). Expect: ``why one slot
empty?'', ``weak-memory changes?''

\question{C2}{Write \code{memcpy}. Now \code{memmove} --- what's the difference?}
Overlap $\rightarrow$ direction choice (Chapter 13 \#16--17).

\question{C3}{Reverse a singly linked list in place.}
Three-pointer iterative (Chapter 13 \#35).

\question{C4}{Count set bits better than 32 iterations.}
Kernighan \code{v\&=v-1} (Chapter 13 \#2).

\question{C5}{Parse/validate a telemetry frame: find sync, check CRC, extract
fields. What if a byte is dropped?}
Find-then-validate, resync on failure, shifts/masks not struct cast
(Chapters 8, 13 \#48/\#53/\#60).

\question{C6}{Write a discrete PID step with anti-windup.}
Integral $\times dt$, clamp, conditional integration (Chapter 10).

\begin{scorecard}
\textbf{Rubric (12 pts):} C1 correct \emph{and} explains the no-lock invariant
(3); C2 both + overlap reasoning (2); C3 correct 3-pointer (2); C4 Kernighan
(1); C5 find/validate/resync \emph{and} no struct-cast (2); C6 correct with
anti-windup (2). Code must be syntactically valid and the candidate must
\emph{dry-run} it on an example. \textbf{Pass $\ge$9.}
\end{scorecard}

\textbf{Common mistakes:} off-by-one / full-vs-empty in the ring buffer; using
\code{char} instead of \code{unsigned char} in \code{memcpy}; losing the
\code{next} pointer in the list reverse (save it first); \code{1<<31} signed UB
in bit code; forgetting to NUL-terminate; not dry-running the code. \emph{Talk
while you write} --- silence reads as being stuck.

%=====================================================================
\section{Mock 5 --- Managerial / Behavioural Round}
%=====================================================================
\emph{~30--45 min, hiring manager. Checks ownership, teamwork, communication,
and fit. Use STAR (Situation, Task, Action, Result). The specific stories come
from INTAKE.md / Chapter 16 --- here, rehearse the structure.}

\subsection*{Question script (answer in STAR)}
\question{M1}{Tell me about a difficult bug you found and fixed.}
STAR; ideally a timing/shared-data or board-level defect with a measured result
(Chapters 4, 11). \textbf{[Real defect from INTAKE.md.]}

\question{M2}{A time you disagreed with a colleague / handled conflict.}
STAR; focus on the resolution and the professional outcome.

\question{M3}{A time you were under deadline pressure. How did you prioritise?}
STAR; show judgement (what you cut, what you protected --- e.g.\ safety/test
coverage).

\question{M4}{Something you learned recently / how you keep current.}
Concrete: a protocol, a tool, this very interview prep --- show growth mindset.

\question{M5}{Why this company? What do you know about what we build?}
\textbf{[Company one-pager from Chapter 16 --- TASL/BEL/HAL/Honeywell/Collins/
Safran/Boeing/Airbus/Bosch/Qualcomm.]}

\question{M6}{Contract-to-permanent / third-party payroll questions.}
\textbf{[From INTAKE.md / Chapter 16 --- factual and positive.]}

\question{M7}{Do you have questions for us?}
Always yes --- ask about the team, the product, the tools, growth. Never ``no.''

\begin{scorecard}
\textbf{Rubric (10 pts):} M1--M3 each a complete STAR with a concrete
\emph{Result} (2 each = 6); M4 shows growth (1); M5 shows genuine company
knowledge (2); M7 has $\ge$2 thoughtful questions ready (1). \textbf{Pass
$\ge$8.}
\end{scorecard}

\textbf{Common mistakes:} stories with no measurable \emph{Result}; blaming
others in the conflict question; ``I have no questions''; not researching the
company; turning ``weakness/why leaving'' negative. Behavioural rounds reward
\emph{structured, positive, specific} --- the same clarity as the technical
ones.

%=====================================================================
\section{Scoring \& the 12-Week On-Ramp}
%=====================================================================
\begin{partnote}
\textbf{Track your mock scores} (target $\ge$80\% each), and re-sit any round you
fail after revisiting its chapters. Per \code{PLAN.md}'s study plan: Weeks 1--3
foundations (Ch 1--3), 4--5 peripherals/serial (4--6), 6--7 CAN/avionics/RTOS
(7--9), 8 control/debug/V\&V (10--12), 9--10 the Coding Gym daily + Python
(13--14), 11 project \& behavioural narratives memorised cold (15--16, once
INTAKE.md is filled), 12 all five mocks here plus redo every failed scorecard.
Speak Tier-2 answers aloud; attempt every coding problem on paper before reading
the solution. Readiness climbs 10\%\,$\rightarrow$\,100\% across the twelve
weeks.
\end{partnote}

\begin{scorecard}
\textbf{Overall readiness gate (sit the real interview when all are true):}
\begin{enumerate}
\item Mock 1 (screen) $\ge$8/10 --- pitch $\le$2\,min, logistics ready.
\item Mock 2 (C \& protocols) $\ge$16/20 --- \code{volatile}-vs-atomic nailed.
\item Mock 3 (domain/projects) $\ge$11/14 --- honest scope boundary clean.
\item Mock 4 (coding) $\ge$9/12 --- ring buffer + frame parse from memory.
\item Mock 5 (behavioural) $\ge$8/10 --- three STAR stories with results.
\item Every chapter scorecard (1--14) passed $\ge$8/10 at least once.
\end{enumerate}
\end{scorecard}

\begin{partnote}
\emph{End of Chapter 17.} The book's technical spine (Chapters 1--14) and mock
execution (17) are complete. \textbf{Chapters 15 (Her Project Deep-Dives) and 16
(Behavioural \& Company Rounds) remain to be written from
\code{inputs/INTAKE.md}} --- they are intentionally gated so every project and
biographical claim is true and defensible. Fill in INTAKE.md's REQUIRED fields,
and those two chapters complete the book.
\end{partnote}
